\subsection*{最初に必要な用語}
\addcontentsline{toc}{subsection}{最初に必要な用語}

\par\smallskip\noindent\refstepcounter{table}\label{tab:ch03-01}\textbf{表 \thetable: 第03章 ソフトウェア設計 表01}\par\smallskip
\begin{tabularx}{\linewidth}{L{34mm}Y}
\toprule
\textbf{用語} & \textbf{初学者向けの説明} \\
\midrule
ソフトウェア設計 & 要求、制約、品質目標を、実装できる構造・振る舞い・インタフェースへ変換する活動である。 \\
設計記述 & 設計結果を記録した文書、モデル、図、表、プロトタイプなどである。英語では Software Design Description、略して SDD と呼ぶ。 \\
構成要素 & システムを構成する部品である。部品は、モジュール、クラス、サービス、コンポーネントなどの形を取る。 \\
インタフェース & ある構成要素が外部とやり取りする接点である。入力、出力、呼び出し方、データ形式、例外条件などを含む。 \\
抽象化 & 注目したい情報だけを取り出し、不要な詳細を隠す考え方である。複雑さを下げるために使う。 \\
モジュール化 & 大きなソフトウェアを、名前を持つ小さな部品に分けることである。理解、分担、保守を容易にする。 \\
カプセル化 & 部品の内部詳細を隠し、外部からは必要な操作だけを見せることである。情報隠蔽とも呼ぶ。 \\
品質属性 & 性能、信頼性、保守容易性、セキュリティ、使いやすさなど、ソフトウェアの性質である。 \\
トレードオフ & ある性質を高めると、別の性質や費用に悪影響が出る関係である。設計では避けられない。 \\
設計根拠 & なぜその設計判断をしたかという理由である。代替案、前提、制約、採用理由、却下理由を含む。 \\
\bottomrule
\end{tabularx}

\subsubsection*{略語}
\addcontentsline{toc}{subsubsection}{略語}

\par\smallskip\noindent\refstepcounter{table}\label{tab:ch03-02}\textbf{表 \thetable: 第03章 ソフトウェア設計 表02}\par\smallskip
\begin{tabularx}{\linewidth}{L{18mm}L{43mm}Y}
\toprule
\textbf{略語} & \textbf{日本語} & \textbf{説明} \\
\midrule
API & アプリケーションプログラミングインタフェース & ソフトウェア部品やサービスを外部から利用するための呼び出し規約である。 \\
AOD & アスペクト指向設計 & 横断的関心事をアスペクトとして扱う設計方法である。 \\
CBD & 構成要素ベース設計 & 独立した構成要素を組み合わせてシステムを作る方法である。 \\
CRC & クラス・責務・協調カード & クラス名、責務、協調相手をカード形式で整理する方法である。 \\
DFD & データフロー図 & データの流れと処理を表す図である。 \\
DSL & 領域固有言語 & 特定領域の概念を直接表すための専用言語である。 \\
ERD & 実体関連図 & データの実体と関係を表す図である。 \\
FOSS & 自由・オープンソースソフトウェア & ソースコード利用、変更、再配布が許されるソフトウェアである。 \\
IDL & インタフェース記述言語 & 構成要素のインタフェースを記述する言語である。 \\
MBD & モデルに基づく設計 & モデルを設計記録の中心に置く方法である。 \\
MDD & モデル駆動開発 & モデルを主要成果物とし、実装やテストへつなげる開発方法である。 \\
OO & オブジェクト指向 & データと操作を持つオブジェクトを中心に考える方法である。 \\
PDL & プログラム設計言語 & プログラムに近い擬似言語で処理を表す記法である。 \\
SDD & ソフトウェア設計記述 & 設計情報を伝達・分析・実装に使うための記録である。 \\
SoC & 関心事の分離 & 複数の観点を分けて考え、複雑さを下げる原則である。 \\
UML & 統一モデリング言語 & 構造や振る舞いを表すための代表的なモデリング記法群である。 \\
\bottomrule
\end{tabularx}
